\documentclass[11pt]{article}
\usepackage[a4paper,margin=27mm]{geometry}
\usepackage{amsmath,amssymb,amsthm,booktabs}
\newtheorem{theorem}{Theorem}
\newtheorem{corollary}[theorem]{Corollary}
\newcommand{\Aut}{\operatorname{Aut}}
\title{The order-four row for O'Nan's group}
\author{}\date{}
\begin{document}\maketitle
Let $G$ be O'Nan's simple group and $A=\Aut(G)$. We count generating
pairs $(R,S)$ under $A$, with $|R|=4$, $|S|=q$ and $|RS|=2$.
The conventions are those of Connor and Leemans, \emph{Algorithmic
enumeration of regular maps}, Ars Math. Contemp. 10 (2016), 211--222.
There is no further quotient by duality, including at equal orders.
An inner or outer pair admits, respectively, an inner or outer
automorphism inverting both rotations. These possibilities are disjoint:
an inverter is unique because its ratio with a second inverter fixes
the generators of $G$.

\begin{theorem}
For $q\ge4$, the complete counts are as follows. The columns give all
pairs, pairs with an inner inverter, pairs with an outer inverter,
and chiral pairs.
\[
\begin{array}{r|rrrr}
q&\mathrm{all}&\mathrm{inner}&\mathrm{outer}&\mathrm{chiral}\\\hline
4&0&0&0&0\\
5&18&0&2&16\\
6&43&0&1&42\\
7&102&1&11&90\\
8&284&6&0&278\\
10&285&5&10&270\\
11&503&15&12&476\\
12&120&5&5&110\\
14&166&3&11&152\\
15&234&12&4&218\\
16&1292&36&0&1256\\
19&846&20&34&792\\
20&554&0&20&534\\
28&370&11&21&338\\
31&359&0&25&334\\
\hline
\mathrm{sum}&5176&114&156&4906
\end{array}
\]
\end{theorem}

\begin{proof}
Use the established correspondence between inversion and generating
triples $(x,y,c)$ of involutions, with $xc=cx$, $x\ne c$, and all three
entries in the same coset of $G$. The rotations are $xy$ and $yc$.
The complete accepted order-four pair partition has three inner
classes, with full centralizers in $A$ of orders $16,144,16$, and two
outer classes, with full centralizers of orders $8,120$.
These are the five actual classes, not merely their abstract types.

For each fixed pair, conjugate its first entry to the acquired fixed
inner or outer involution. The complete commuting domains contain,
respectively, $1750$ and $1463$ eligible third involutions. Partition
the domain under the full pair centralizer $H$. The five complete
partitions contain, in the same order, $126,27,134,210,22$ orbits.
The stabilizer acts freely on every generating triple, so the number
of generating classes of second order $q$ is the number of such
third involutions divided by $|H|$.

Every orbit representative is supplied by a word in the acquired
permutation action of degree $245520$. The reader checks the involution
and commutation identities and both rotation orders. It also computes
the orbit of a fixed point under the two rotations. The socle orbit
has size $122760$, and none of the thirteen acquired maximal-subgroup
orders is divisible by $122760$. Thus the rotations generate $G$
if and only if their orbit has that size: a proper generated subgroup
would lie in a maximal subgroup whose order must be divisible by its
point-orbit size. This is a complete generation criterion, with the
maximal roster an explicit borrowed input.

The resulting numbers of generating classes in the five partitions
are $48,0,66,150,6$. Each generating orbit is checked individually to
have size $|H|$. Adding its contribution at its exact second order
gives the inner and outer columns. The zero contribution of the
centralizer of order $144$ is reused from its accepted certificate;
the other four journals are reconstructed literally. All domain
orbits are included. Finally subtract the two disjoint reflexible
counts from the previously proved total in each cell. This gives the
chiral column and the displayed sums.
\end{proof}

\begin{corollary}
There are no inner-reflexible pairs of types $\{4,5\}$ or $\{4,6\}$.
Their outer and chiral counts are $(2,16)$ and $(1,42)$, respectively.
In particular type $\{4,5\}$ admits a chiral map.
\end{corollary}

These statements recover the corresponding assertions in Section 5
of the source and its complete order-four totals in Section 7.
The source's historical timing is attributed, not reproduced. The
new certificate uses the shared centralizer domains and the accepted
fixed-pair theorem; the table of total counts, finite group inputs,
and inversion correspondence are imported with their existing
attribution. No novelty claim is made for the underlying correspondence.

Combining the two new inner zeros with the previously established
$\{3,10\}$ zero and the accepted exclusions involving orders $20$ or
$31$ leaves at most $95$ possible inner-reflexible types. This is an
upper bound only. The claimed existence of every one of those types
still needs its own witnesses. The analogous $98$-type outer spectrum,
the order-five row and the remaining whole-source claims are also
separate obligations. The result here grants no whole-original,
open-conjecture or Lean credit.
\end{document}
